\documentclass[a4paper,12pt]{article}

% ---------- Paquetes ----------
\usepackage[utf8]{inputenc}
\usepackage[T1]{fontenc}
\usepackage[spanish]{babel}
\usepackage{geometry}
\geometry{top=2.5cm, bottom=2.5cm, left=2.5cm, right=2.5cm}
\usepackage{hyperref}
\hypersetup{
    colorlinks=true,
    linkcolor=blue,
    urlcolor=blue,
    pdftitle={Manual del Proyecto - Operador de Marketing e Inteligencia de Inventario D2C},
    pdfauthor={Sistema}
}
\usepackage{listings}
\usepackage{xcolor}
\usepackage{graphicx}
\usepackage{titlesec}
\usepackage{fancyhdr}
\usepackage{enumitem}
\usepackage{verbatim}

% ---------- Configuración listings ----------
\lstset{
    basicstyle=\ttfamily\small,
    breaklines=true,
    frame=single,
    numbers=left,
    numberstyle=\tiny,
    backgroundcolor=\color{gray!10},
    rulecolor=\color{gray!30},
    keywordstyle=\color{blue},
    commentstyle=\color{green!50!black},
    stringstyle=\color{red!60!black}
}

% ---------- Encabezado / Pie ----------
\pagestyle{fancy}
\fancyhf{}
\fancyhead[L]{\leftmark}
\fancyhead[R]{\thepage}
\renewcommand{\headrulewidth}{0.4pt}

% ---------- Portada ----------
\title{\Huge\textbf{Operador de Marketing e\\Inteligencia de Inventario D2C}\\[0.8cm]
\Large Manual del Proyecto\\[0.3cm]
\large Automatización de Marketing y Gestión de Inventario para Retail D2C}
\author{Sistema de Agente IA --- Arquitectura de 3 Capas}
\date{\today}

\begin{document}

\maketitle
\thispagestyle{empty}
\newpage

\tableofcontents
\newpage

% ============================================================
\section{Introducción}
% ============================================================

Este proyecto implementa un \textbf{agente de inteligencia artificial} especializado en la automatización de marketing y la gestión de inventario para una tienda \textbf{Direct-to-Consumer (D2C)} de ropa de golf. La tienda opera simultáneamente en \textbf{Shopify} (tienda online), \textbf{Whatnot} (ventas en vivo) y \textbf{Meta Ads} (publicidad).

El objetivo central es \textbf{reducir las tareas manuales} del equipo humano y optimizar la precisión de las operaciones y el marketing mediante un sistema de alertas, borradores automáticos y sincronización de datos, manteniendo siempre la supervisión humana como paso obligatorio antes de cualquier acción ejecutable.

\subsection{Propósito del Sistema}

\begin{itemize}
    \item Actuar como puente entre la intención del usuario y la ejecución técnica.
    \item Separar la lógica probabilística (IA) de la ejecución determinista (scripts).
    \item Proveer un marco reproducible, con estado trazable y capacidad de autocorrección.
\end{itemize}

\subsection{Stack Tecnológico}

\begin{itemize}
    \item \textbf{Lenguaje:} Python 3.12+
    \item \textbf{Librerías principales:} ShopifyAPI, python-dotenv, psutil.
    \item \textbf{Infraestructura:} Docker (sandbox para compilación nativa), Conda (gestión de entornos).
    \item \textbf{Automatización:} Compatible con Zapier / Make.com.
    \item \textbf{Memoria persistente:} ChromaDB.
    \item \textbf{SO base:} Linux (Kernel compatible con Linux Mint).
\end{itemize}

% ============================================================
\section{Estructura del Proyecto}
% ============================================================

La siguiente es la estructura de directorios del proyecto:

\begin{lstlisting}
.
+-- .agent/                      # Instrucciones de sistema y contexto del agente
|   +-- AGENT_FRAMEWORK.md       # Arquitectura de 3 capas y setup
|   +-- AGENT_INSTRUCTIONS.md    # Protocolo de agente y memoria evolutiva
|   +-- Dockerfile               # Imagen para sandbox Docker
+-- .env                         # Credenciales y configuracion del agente
+-- .gitignore
+-- directives/                  # SOPs en YAML (Layer 1)
|   +-- check_shopify_inventory.yaml
+-- docs/
|   +-- OBJETIVO.md              # Prompt principal del agente
+-- execution/                   # Scripts deterministas en Python (Layer 3)
|   +-- env_diagnostic.py
|   +-- shopify_get_inventory.py
+-- requirements.txt             # Dependencias del proyecto
+-- git-update.sh                # Script de actualizacion Git
+-- update_repo.sh               # Script de sincronizacion de repositorio
+-- env_diagnostic.py            # Diagnostico de entorno
+-- README.md
+-- manual.tex                   # Este documento
\end{lstlisting}

\subsection{Descripción de Directorios}

\begin{description}
    \item[.agent/] Contiene las instrucciones del sistema y la definición del marco de trabajo del agente (arquitectura de 3 capas, protocolos de ejecución, memoria).
    \item[directives/] Archivos YAML que actúan como Procedimientos Operativos Estandarizados (SOP). Cada flujo de trabajo repetible tiene su propia directiva.
    \item[execution/] Scripts Python deterministas con una única responsabilidad. Son invocados por la capa de orquestación.
    \item[.tmp/] Archivos temporales y estado de ejecución (regenerables).
    \item[docs/] Documentación del proyecto.
\end{description}

% ============================================================
\section{Arquitectura del Sistema}
% ============================================================

El sistema se basa en una \textbf{arquitectura de 3 capas} diseñada para maximizar la fiabilidad, separando la lógica probabilística del LLM de la ejecución determinista.

\subsection{Capa 1: Directivas (directives/) --- \textit{Qué hacer}}

Son SOPs en formato YAML, con un archivo por flujo de trabajo repetible. Cada directiva contiene:

\begin{itemize}
    \item \textbf{Goal:} Objetivo principal del flujo.
    \item \textbf{Required inputs:} Lista de datos necesarios (nombre y descripción).
    \item \textbf{Steps:} Secuencia de pasos que invocan scripts de execution/.
    \item \textbf{Expected outputs:} Resultados esperados.
    \item \textbf{Edge cases:} Casos límite y protocolos de recuperación.
\end{itemize}

Ejemplo de estructura de una directiva:

\begin{lstlisting}
goal: "Obtener niveles de stock actuales de Shopify"
required_inputs:
  - name: "inventory_threshold"
    description: "Stock minimo antes de alertar (ej. 5)"
steps:
  - step: 1
    description: "Consultar API de Shopify"
    script: "execution/shopify_get_inventory.py"
expected_outputs:
  - "Archivo JSON en .tmp/inventory_report.json"
edge_cases:
  - case: "Error de conexion"
    protocol: "Reintentar hasta 3 veces"
\end{lstlisting}

\subsection{Capa 2: Orquestación (El Agente) --- \textit{Toma de decisiones}}

Esta capa es el \textbf{orquestador inteligente}. Sus funciones incluyen:

\begin{itemize}
    \item Leer la directiva apropiada para la tarea solicitada.
    \item Elegir las herramientas de execution/.
    \item Ejecutar flujos multietapa.
    \item Validar entradas y salidas.
    \item Gestionar errores y autocorrección (self-annealing).
    \item Guardar estado en \texttt{.tmp/run\_state.json}.
    \item Actualizar directivas con conocimiento aprendido.
\end{itemize}

\textbf{Regla fundamental:} No inventar lógica que deba vivir en los scripts. El agente conecta intención e implementación.

\subsection{Capa 3: Ejecución (execution/) --- \textit{Hacer el trabajo}}

Scripts Python deterministas con los siguientes requisitos:

\begin{itemize}
    \item Entradas por argumentos de línea de comandos.
    \item Secretos en \texttt{.env} (nunca hardcodeados).
    \item Salidas por stdout (preferiblemente en JSON).
    \item Códigos de salida: 0 = éxito, 1+ = fallos categorizados.
    \item Validación propia de salidas.
    \item Fallo ruidoso ante condiciones anómalas.
\end{itemize}

% ============================================================
\section{Configuración Inicial}
% ============================================================

\subsection{Requisitos Previos}

\begin{itemize}
    \item Python 3.12 o superior.
    \item Conda (recomendado) o virtualenv.
    \item Docker (para sandbox de compilación).
    \item Acceso a Internet y a GitHub.
\end{itemize}

\subsection{Entorno Virtual}

\begin{lstlisting}[language=bash]
# Crear entorno Conda
conda create --name d2c_env python=3.12 -y

# Activar entorno
conda activate d2c_env

# Instalar dependencias
pip install -r requirements.txt
\end{lstlisting}

\subsection{Docker Sandbox}

Para ejecución de código C/C++ de forma aislada:

\begin{lstlisting}[language=bash]
docker build -t pcb_sandbox .agent/
\end{lstlisting}

\subsection{Archivo .env}

El archivo \texttt{.env} contiene las credenciales necesarias:

\begin{itemize}
    \item \texttt{SHOPIFY\_SHOP\_URL} --- URL de la tienda Shopify.
    \item \texttt{SHOPIFY\_ACCESS\_TOKEN} --- Token de acceso a la API de Shopify.
    \item Claves de API para LLMs (Google, OpenAI, Anthropic, etc.).
    \item Configuración de Telegram (bot token, chat ID).
    \item Configuración de gestión de contexto (límites de mensajes y caracteres).
\end{itemize}

% ============================================================
\section{Módulos y Funcionalidades}
% ============================================================

\subsection{Gestión de Inventario Shopify}

El script \texttt{execution/shopify\_get\_inventory.py} consulta la API de Shopify y genera un reporte de inventario.

\subsubsection{Uso}

\begin{lstlisting}[language=bash]
python execution/shopify_get_inventory.py --threshold 5
\end{lstlisting}

\subsubsection{Parámetros}

\begin{description}
    \item[--threshold] Nivel de stock por debajo del cual se considera ``bajo inventario'' (por defecto: 5).
\end{description}

\subsubsection{Salida}

Genera un archivo JSON en \texttt{.tmp/inventory\_report.json} con la siguiente estructura:

\begin{lstlisting}
{
  "out_of_stock": [
    { "id": 123, "title": "Producto - Variante",
      "sku": "GOLF-001", "stock": 0 }
  ],
  "low_stock": [
    { "id": 456, "title": "Producto2 - Variante2",
      "sku": "GOLF-002", "stock": 3 }
  ],
  "summary": { "total_checked": 150 }
}
\end{lstlisting}

\subsubsection{Manejo de Errores}

\begin{itemize}
    \item \textbf{Error de conexión:} Reintentar hasta 3 veces.
    \item \textbf{API Rate Limit:} Detectar header \texttt{Retry-After} y pausar automáticamente.
    \item \textbf{Credenciales faltantes:} Mensaje de error claro en JSON.
\end{itemize}

\subsection{Alertas de Inventario y Publicidad}

El sistema cruza el gasto en anuncios de Meta Ads contra el stock disponible en Shopify para:

\begin{itemize}
    \item Identificar anuncios activos de productos agotados o con bajo stock.
    \item Generar \textbf{Alertas Prioritarias} cuando hay discrepancia entre gasto publicitario y stock.
    \item Detectar productos recién repuestos para activar campañas de marketing.
\end{itemize}

\subsection{Generación de Contenido de Marketing}

Basado en la disponibilidad de inventario, el agente genera borradores de:

\begin{itemize}
    \item Campañas de Email (flujos Klaviyo).
    \item Anuncios para Meta (Facebook/Instagram).
    \item Contenido para redes sociales.
    \item SMS marketing.
\end{itemize}

Todo el contenido se alinea con la identidad de la marca de ropa de golf y está diseñado para ser atractivo para la comunidad del golf.

\subsection{Asistente de Soporte al Cliente}

El agente redacta borradores de respuesta a consultas de soporte extrayendo información de:

\begin{itemize}
    \item Pedidos de Shopify (estado, productos, fechas).
    \item Políticas de la empresa (devoluciones, envíos).
    \item Historial de respuestas anteriores.
\end{itemize}

El usuario solo debe revisar y enviar.

\subsection{Sincronización Whatnot a Shopify}

El sistema ayuda a estructurar la lógica para:

\begin{itemize}
    \item Mapear ventas realizadas en vivo en Whatnot hacia pedidos en Shopify.
    \item Asegurar que los datos de inventario estén listos para imprimir hojas de stock para shows en vivo.
\end{itemize}

% ============================================================
\section{Protocolo de Diagnóstico de Entorno}
% ============================================================

Antes de cualquier ejecución, el agente debe conocer el entorno real del sistema. El protocolo es:

\begin{enumerate}
    \item Generar \texttt{env\_diagnostic.py} que recolecte: SO, arquitectura, versión de Python, entorno Conda, paquetes críticos, RAM, CPU, GPU, permisos, conectividad.
    \item El usuario ejecuta el script y pega la salida.
    \item El agente ajusta su ejecución a las versiones y limitaciones confirmadas.
    \item Si cambia el contexto de trabajo o surge un \texttt{ImportError}, solicitar nuevo diagnóstico.
\end{enumerate}

% ============================================================
\section{Memoria Persistente (ChromaDB)}
% ============================================================

El agente utiliza ChromaDB para memoria persistente:

\begin{itemize}
    \item \textbf{Consulta inicial:} Antes de proponer una solución, consulta la memoria para experiencias pasadas.
    \item \textbf{Registro de aprendizaje:} Errores críticos y limitaciones técnicas se guardan con \texttt{save\_memory.yaml}.
    \item \textbf{Autocorrección:} Si un script falla, busca fallos similares antes de intentar una nueva solución.
\end{itemize}

Directivas disponibles:

\begin{itemize}
    \item \texttt{save\_memory.yaml} --- Guardar aprendizaje en memoria.
    \item \texttt{query\_memory.yaml} --- Recuperar información relevante.
    \item \texttt{list\_memories.yaml} --- Listar recuerdos recientes.
    \item \texttt{delete\_memory.yaml} --- Eliminar un recuerdo por ID.
\end{itemize}

% ============================================================
\section{Formato de Salida}
% ============================================================

\subsection{Borradores para Revisión Humana}

Todo contenido generado por el agente (emails, publicaciones, respuestas de soporte) se presenta como \textbf{borrador}. El usuario mantiene el control final y debe aprobar antes de ejecutar.

\subsection{Alertas Prioritarias}

Cuando se detecta una discrepancia entre el gasto publicitario y el stock disponible, se genera una \textbf{Alerta Prioritaria} que se resalta en la salida para acción inmediata.

\subsection{Estado de Ejecución}

Cada paso de ejecución se registra en \texttt{.tmp/run\_state.json} con:

\begin{itemize}
    \item Timestamp.
    \item Script ejecutado.
    \item Código de salida (exit\_code).
\end{itemize}

% ============================================================
\section{Mantenimiento y Git}
% ============================================================

\subsection{Scripts de Actualización}

El proyecto incluye dos scripts para la gestión de versiones:

\subsubsection{git-update.sh}

Script principal que:
\begin{enumerate}
    \item Detecta la rama actual.
    \item Commitea cambios locales con mensaje ``WIP: guardar cambios antes de pull''.
    \item Delega en \texttt{update\_repo.sh}.
\end{enumerate}

\subsubsection{update\_repo.sh}

Script de sincronización con soporte para:

\begin{lstlisting}[language=bash]
./update_repo.sh                    # Pull + commit + push
./update_repo.sh --no-pull --push   # Solo commit y push
./update_repo.sh --dry-run          # Vista previa sin ejecutar
./update_repo.sh --confirm          # Modo interactivo con confirmacion
./update_repo.sh -m "Mensaje"       # Mensaje de commit personalizado
\end{lstlisting}

\textbf{Parámetros:}

\begin{description}
    \item[-{}-dir DIR] Directorio del repositorio (por defecto: actual).
    \item[-{}-remote REMOTE] Remote Git (por defecto: origin).
    \item[-{}-branch BRANCH] Rama (por defecto: actual).
    \item[-{}-push] Hacer push después del commit.
    \item[-{}-no-pull] Omitir pull inicial.
    \item[-{}-dry-run] Vista previa sin cambios.
    \item[-{}-confirm] Modo interactivo.
\end{description}

% ============================================================
\section{Principios de Operación}
% ============================================================

\subsection{Autocorrección (Self-Annealing)}

\begin{itemize}
    \item Presupuesto máximo de 3 reintentos por tarea.
    \item Clasificar fallos en: \textbf{Lógica}, \textbf{Entorno} o \textbf{Recursos}.
    \item Analizar causa raíz antes de corregir.
    \item Fiabilidad sobre velocidad: mejor preguntar que continuar con datos inconsistentes.
\end{itemize}

\subsection{Reutilización}

Antes de escribir un nuevo script, verificar si existe en \texttt{execution/} una herramienta que pueda reutilizarse.

\subsection{Actualización de Directivas}

Cada hallazgo (errores comunes, límites, mejoras) debe registrarse en las directivas sin sobrescribir versiones anteriores.

\subsection{Validación Continua}

Después de cada paso, confirmar:
\begin{itemize}
    \item Esquema correcto.
    \item Cantidades coherentes.
    \item Archivos en las rutas esperadas.
    \item Tiempos y fechas razonables.
\end{itemize}

% ============================================================
\section{Seguridad}
% ============================================================

\begin{itemize}
    \item \textbf{Sanitización de entradas:} Verificar que rutas y parámetros no contengan caracteres de escape maliciosos (\texttt{;}, \texttt{\&}, \texttt{|}, etc.).
    \item \textbf{Validación de tipos:} Los scripts deben forzar tipos de datos (Type Hinting) para evitar errores de casting en runtime.
    \item \textbf{Credenciales:} Nunca hardcodear secretos en los scripts; usar \texttt{.env}.
    \item \textbf{Aislamiento:} Uso de Docker para compilación de código no confiable.
\end{itemize}

% ============================================================
\section{Apéndice: Dependencias}
% ============================================================

El archivo \texttt{requirements.txt} contiene:

\begin{lstlisting}
psutil
python-dotenv
ShopifyAPI
\end{lstlisting}

% ============================================================
\section{Apéndice: Referencia Rápida de Comandos}
% ============================================================

\begin{center}
\begin{tabular}{|l|l|}
\hline
\textbf{Comando} & \textbf{Descripción} \\
\hline
\texttt{conda create -n d2c\_env python=3.12 -y} & Crear entorno \\
\texttt{conda activate d2c\_env} & Activar entorno \\
\texttt{pip install -r requirements.txt} & Instalar dependencias \\
\texttt{python execution/shopify\_get\_inventory.py --threshold 5} & Consultar inventario \\
\texttt{./git-update.sh} & Actualizar repositorio \\
\texttt{./update\_repo.sh --dry-run} & Vista previa de actualización \\
\texttt{docker build -t pcb\_sandbox .agent/} & Construir sandbox Docker \\
\hline
\end{tabular}
\end{center}

\end{document}
